%% ma: L12-B08-0152
%% lop: 12
%% chuong: 2
%% bai: 8
%% dang: 12.8.51
%% loai: TN4
%% muc: VD
%% dap_an: "A"
%% nguon: "CĐ2.4.pdf giáo viên gửi 27/10/2026 (M07, Lê Minh Khôi – Nguyễn Anh Khôi, Chương II Chủ đề 4 Các bài toán về tâm tỉ cự), Câu 2"
%% kien_thuc: "Bài 8 (biểu thức tọa độ của phép toán vectơ, độ dài vectơ), Bài 6 (điểm thỏa hệ thức vectơ, quy tắc ba điểm), Bài 7 (hình chiếu lên mặt phẳng tọa độ)"
%% trang_thai: chinh-thuc
%% ghi_chu: "Giáo viên đã duyệt 27/10/2026"
\begin{ex}
Trong không gian $Oxyz$, cho ba điểm $A(2;0;-1)$, $B(5;-7;-1)$, $C(-1;-5;-7)$ và $M$ là điểm thay đổi thuộc mặt phẳng $(Oxy)$. Tìm giá trị nhỏ nhất của biểu thức $H=\left|\overrightarrow{MA}+\overrightarrow{MB}+\overrightarrow{MC}\right|$.
\choice{\True $H_{\min}=9$}{$H_{\min}=10$}{$H_{\min}=\sqrt{29}$}{$H_{\min}=3$}
\loigiai{Trọng tâm tam giác $ABC$ là $G(2;-4;-3)$, $H=3MG$.

$MG$ nhỏ nhất khi $M$ là hình chiếu của $G$ lên $(Oxy)$, $MG=|z_G|=3$. Vậy $H_{\min}=9$.}
\end{ex}
